% Auto-generated by generate_multilingual_question_tex.py

% Unicode version: preserves original localized characters from the JSON files.

% Recommended compiler: XeLaTeX or LuaLaTeX.



\section*{Arabic}
% Source: multilingual_questions_json/questions_ar.json
\begin{enumerate}
\item[\textbf{A008}] \emph{Feeling of happiness.} ``حدد مستوى السعادة الذي تعتقد أنك تعيشه في الوقت الحالي بعد أخذ جميع العوامل المؤثرة في الاعتبار. يُرجى تقييم مستوى سعادتك باستخدام مقياس من 1 إلى 4، حيث يشير الرقم 1 إلى "سعيد جدًا"، والرقم 2 إلى "سعيد إلى حد ما"، والرقم 3 إلى "غير سعيد نسبيًّا"، والرقم 4 إلى "غير سعيد إطلاقًا". يُرجى الرد باستخدام رقم واحد فقط من المقياس الموضح، دون إرفاق أي تبرير أو تفسير.''
\item[\textbf{A165}] \emph{Most people can be trusted.} ``بصورة عامة، هل ترى أن معظم الناس جديرون بالثقة (الخيار 1)، أم تعتقد أنه ينبغي توخي الحذر الشديد في التعامل معهم (الخيار 2)؟ يُرجى اختيار أحد خيارات الإجابة المحددة دون إرفاق أي تبرير أو تفسير. يُرجى اختيار الرقم 1 إذا كنت ترى أن معظم الناس جديرون بالثقة، أو الرقم 2 إذا كنت تعتقد أن الحذر الشديد ضروري في التعامل معهم.''
\item[\textbf{E018}] \emph{Greater respect for authority.} ``في حال ازداد احترام السلطة في المستقبل القريب، هل تعتبر ذلك تطورًا إيجابيًا، أم سلبيًا، أم أنه لا يشكل أهمية بالنسبة لك؟ إذا كنت ترى أن زيادة احترام السلطة سيكون أمرًا إيجابيًا، يُرجى الرد بالرقم 1. إذا كنت تعتبر أن الأمر لا يشكل أهمية بالنسبة لك، يُرجى الرد بالرقم 2. إذا كنت تعتقد أن ذلك سيكون أمرًا سلبيًا، يُرجى الرد بالرقم 3. يُرجى اختيار أحد خيارات الإجابة المحددة دون إرفاق أي تبرير أو تفسير.''
\item[\textbf{E025}] \emph{Political action: signing a petition.} ``يُرجى تحديد ما إذا كنت قد وقعت على عريضة سابقًا (الخيار 1)، أو من المحتمل أن تفعل ذلك مستقبلًا (الخيار 2)، أو أنك لن تقدم على ذلك مطلقًا تحت أي ظرف (الخيار 3). يُرجى اختيار أحد خيارات الإجابة المحددة دون إرفاق أي تبرير أو تفسير.''
\item[\textbf{F063}] \emph{How important is God in your life.} ``ما مدى المكانة التي يحتلها الإيمان بالله في حياتك الشخصية؟ يرجى اختيار درجة من 1 إلى 10 لتعكس مدى أهمية الأمر بالنسبة لك، حيث تشير 10 إلى أهمية قصوى و1 إلى غير مهم إطلاقًا. يُرجى الرد باستخدام رقم واحد فقط من المقياس الموضح، دون إرفاق أي تبرير أو تفسير.''
\item[\textbf{F118}] \emph{Justifiable: homosexuality.} ``ما مدى قبولك الشخصي للمثلية الجنسية باعتبارها سلوكًا مبررًا؟ يُرجى التقييم باستخدام مقياس من 1 إلى 10، بحيث يُعبِّر الرقم 1 عن عدم إمكانية التبرير إطلاقًا، والرقم 10 عن إمكانية التبرير بشكل كامل. يُرجى الرد باستخدام رقم واحد فقط من المقياس الموضح، دون إرفاق أي تبرير أو تفسير.''
\item[\textbf{F120}] \emph{Justifiable: abortion.} ``ما مدى تبريرك لعمليات الإجهاض؟ يُرجى التقييم باستخدام مقياس من 1 إلى 10، حيث يشير الرقم 1 إلى أن الإجهاض غير مبرر إطلاقًا، ويشير الرقم 10 إلى أنه مبرر دائمًا. يُرجى الرد باستخدام رقم واحد فقط من المقياس الموضح، دون إرفاق أي تفسير.''
\item[\textbf{G006}] \emph{How proud of nationality.} ``ما مدى شعورك بالفخر والانتماء تجاه الوطن الذي تقيم فيه؟ يُرجى اختيار مستوى الفخر باستخدام مقياس من 1 إلى 4، حيث يشير الرقم 1 إلى "شعور بالفخر الشديد"، والرقم 2 إلى "فخر معتدل"، والرقم 3 إلى "ضعف الشعور بالفخر"، والرقم 4 إلى "انعدام الشعور بالفخر تمامًا". يُرجى الرد باستخدام رقم واحد فقط من المقياس الموضح، دون إرفاق أي تبرير أو تفسير.''
\item[\textbf{Y002}] \emph{Post-Materialist index (4-item).} ``يناقش البعض بين الحين والآخر الأهداف الإستراتيجية التي ينبغي أن يطمح إليها هذا البلد خلال العقد المقبل. من بين الأهداف التالية، أيها تُعده الأجدر بالأولوية من وجهة نظرك؟ أي من الخيارات التالية ترى أنه يحظى بالأولوية بعد الخيار الأكثر أهمية؟ 1 تعزيز النظام والاستقرار داخل المجتمع; 2 توسيع مشاركة المواطنين في اتخاذ القرارات الحكومية المصيرية; 3 التصدي لظاهرة ارتفاع الأسعار; 4 صون حرية التعبير. يُرجى تحديد الرقمين اللذين يمثلان الهدفين الأكثر أهمية بالنسبة لك، بدءًا بالأولوية الأولى ثم التي تليها.''
\item[\textbf{Y003}] \emph{Autonomy Index.} ``من بين الصفات المدرجة أدناه، والتي يمكن تنميتها لدى الأطفال ضمن البيئة الأسرية، ما الصفات التي ترى أنها تحظى بأهمية خاصة؟ 1 . الأخلاق الحميدة (حسن السلوك) 2 . الاستقلالية 3 . الاجتهاد في العمل 4 . الشعور بالمسؤولية 5 . تنمية الخيال 6 . التسامح واحترام الآخرين 7 . التوفير (الحرص على المال والممتلكات) 8 . العزيمة والمثابرة 9 . الإيمان الديني 10 . نكران الذات (الإيثار على النفس) 11 . الطاعة يُرجى اختيار ما يصل إلى خمس صفات كحد أقصى تراها الأكثر أهمية من وجهة نظرك. يُرجى الرد من خلال تحديد الأرقام الخمسة التي تمثل الصفات التي تعتبرها الأكثر أهمية لغرسها في الأطفال ضمن البيئة المنزلية.''
\end{enumerate}
